\documentclass[10pt,a4paper]{article}
\usepackage[margin=2cm]{geometry}
\usepackage{amsmath,booktabs,array}
\title{Lobster STG AB/PD model with transient LP inhibition}
\author{Brian2 student tutorial based on Prinz, Bucher \& Marder (2004)}
\date{}
\begin{document}\maketitle

\section*{Scope and selected cell}
The biological pacemaker kernel contains one anterior burster (AB) and two
pyloric dilator (PD) neurons. Following the paper, they are lumped into one
single-compartment AB/PD model. The tutorial selects \textbf{AB/PD model 3};
the paper provides five valid AB/PD conductance combinations.

\section*{Membrane and calcium equations}
\begin{align}
C\frac{dV}{dt}&=-\left(I_{Na}+I_{CaT}+I_{CaS}+I_A+I_{KCa}+I_{Kd}
+I_H+I_{leak}+I_{syn}\right),\\
I_i&=\bar g_i m_i^{p_i}h_i(V-E_i),\qquad
I_{leak}=g_{leak}(V-E_{leak}),\\
\tau_{Ca}\frac{d[Ca]}{dt}&=-f(I_{CaT}+I_{CaS})-[Ca]+[Ca]_0,\\
E_{Ca}&=12.2\,\mathrm{mV}\ln\!\left(\frac{[Ca]_{out}}{[Ca]}\right).
\end{align}
The area is $0.628\times10^{-3}$ cm$^2$, $C=0.628$ nF,
$\tau_{Ca}=200$ ms, $f=14.96\,\mu$M/nA, $[Ca]_0=0.05\,\mu$M, and
$[Ca]_{out}=3$ mM. Reversals are $E_{Na}=50$, $E_K=-80$, $E_H=-20$, and
$E_{leak}=-50$ mV. The calcium reversal is dynamic.

\section*{AB/PD model 3 maximal conductances}
Values below are reproduced from Table 2 of the supplied paper.
\begin{center}
\begin{tabular}{lrrrrrrrr}
\toprule
Current & Na & CaT & CaS & A & K(Ca) & Kd & H & leak\\
\midrule
$\bar g$ (mS/cm$^2$) & 200 & 2.5 & 4 & 50 & 5 & 50 & 0.01 & 0\\
Power $p$ & 3 & 3 & 3 & 3 & 4 & 4 & 1 & --\\
Inactivation $h$? & yes & yes & yes & yes & no & no & no & --\\
\bottomrule
\end{tabular}
\end{center}

\section*{Channel gates}
Every gate obeys
\begin{equation}
\tau_x(V,[Ca])\frac{dx}{dt}=x_\infty(V,[Ca])-x.
\end{equation}
Writing $v=V/(1\,\mathrm{mV})$, the steady-state functions are
\begin{align*}
m_{Na,\infty}&=[1+e^{(v+25.5)/-5.29}]^{-1},&
h_{Na,\infty}&=[1+e^{(v+48.9)/5.18}]^{-1},\\
m_{CaT,\infty}&=[1+e^{(v+27.1)/-7.2}]^{-1},&
h_{CaT,\infty}&=[1+e^{(v+32.1)/5.5}]^{-1},\\
m_{CaS,\infty}&=[1+e^{(v+33)/-8.1}]^{-1},&
h_{CaS,\infty}&=[1+e^{(v+60)/6.2}]^{-1},\\
m_{A,\infty}&=[1+e^{(v+27.2)/-8.7}]^{-1},&
h_{A,\infty}&=[1+e^{(v+56.9)/4.9}]^{-1},\\
m_{KCa,\infty}&=\frac{[Ca]}{[Ca]+3\,\mu\mathrm{M}}\times
[1+e^{(v+28.3)/-12.6}]^{-1},\\
m_{Kd,\infty}&=[1+e^{(v+12.3)/-11.8}]^{-1},&
m_{H,\infty}&=[1+e^{(v+75)/5.5}]^{-1}.
\end{align*}
The complete time-constant expressions are written explicitly beside the
Brian2 equations in the accompanying Python file; all are in milliseconds and
come from the 2003 model specification referenced by the supplied paper.

\section*{Transient LP-to-PD glutamatergic inhibition}
LP is glutamatergic, so the paper specifies $E_s=-70$ mV,
$k_-=1/40$ ms$^{-1}$, $V_{th}=-35$ mV, and $\Delta=5$ mV:
\begin{align}
I_{syn}&=g_s s(V-E_s),\\
\frac{ds}{dt}&=\frac{s_\infty(V_{LP})-s}{\tau_s},&
s_\infty(V_{LP})&=\frac{1}{1+e^{(V_{th}-V_{LP})/\Delta}},&
\tau_s&=\frac{1-s_\infty}{k_-}.
\end{align}
The paper tests $g_s\in\{0,3,10,30,100\}$ nS; this lesson uses 30 nS.
The settings dictionary controls when an idealized LP active waveform begins
and ends. During that interval $V_{LP}=20$ mV; otherwise $V_{LP}=-60$ mV.

\section*{Implementation note}
The original database initialized each cell from a saved point on its limit
cycle. This standalone lesson instead performs a 5 s unrecorded settling run.
The reported spike threshold ($-20$ mV) is a measurement convention only;
spikes and bursts arise from the conductance dynamics without voltage reset.

\section*{Sources}
Prinz, A.A., Bucher, D. \& Marder, E. (2004). \emph{Similar network activity
from disparate circuit parameters}. Nature Neuroscience 7, 1345--1352.
The channel kinetics are from Prinz, A.A., Billimoria, C.P. \& Marder, E.
(2003). Journal of Neurophysiology 90, 3998--4015.
\end{document}
